% !TEX program = xelatex
% =====================================================================
% LLM Theory Seminar Week 4: Transformer Expressivity
%
% Compile:
%   xelatex transformer_expressivity_notes.tex
%   xelatex transformer_expressivity_notes.tex
% or
%   latexmk -xelatex transformer_expressivity_notes.tex
%
% Korean typesetting: kotex + XeLaTeX.
%
% Fixed layout reference / inspiration:
%   - Ben Jones, "lecture-notes" (Overleaf, CC BY 4.0)
%   - CMU-Notes/notes-template (structured university course notes)
% This file intentionally keeps the approved Week-1/2/3 seminar-note layout.
% =====================================================================

\documentclass[11pt,a4paper,oneside,openany]{memoir}

\usepackage{kotex}
\usepackage{amsmath,amssymb,amsthm,mathtools,bm}
\usepackage{booktabs,longtable,array,tabularx,makecell}
\newcolumntype{Y}{>{\raggedright\arraybackslash}X}
\usepackage{enumitem}
\usepackage[most]{tcolorbox}
\usepackage[dvipsnames]{xcolor}
\usepackage{xurl}
\usepackage[unicode]{hyperref}
\usepackage{cleveref}
\usepackage{needspace}

% ---------- page geometry ----------
\setlrmarginsandblock{27mm}{27mm}{*}
\setulmarginsandblock{24mm}{28mm}{*}
\checkandfixthelayout

\setlength{\parindent}{0pt}
\setlength{\parskip}{0.48em}
\setlength{\emergencystretch}{2em}
\setlength{\headheight}{15pt}
\linespread{1.055}\selectfont
\setlist[itemize]{topsep=0.28em,itemsep=0.12em,parsep=0pt,leftmargin=1.55em}
\setlist[enumerate]{topsep=0.28em,itemsep=0.16em,parsep=0pt,leftmargin=1.75em}

% ---------- restrained lecture-note palette ----------
\definecolor{NoteBlue}{HTML}{294E6B}
\definecolor{NoteBlueLight}{HTML}{F2F6F9}
\definecolor{NoteGray}{HTML}{5A6268}
\definecolor{NoteGrayLight}{HTML}{F6F6F4}
\definecolor{NoteWarm}{HTML}{7A6545}
\definecolor{NoteWarmLight}{HTML}{FAF7F0}
\definecolor{NoteRule}{HTML}{C8CDD2}

\hypersetup{
  colorlinks=true,
  linkcolor=NoteBlue,
  citecolor=NoteBlue,
  urlcolor=NoteBlue,
  pdftitle={Transformer Expressivity},
  pdfauthor={LLM Theory Seminar}
}

% ---------- running heads ----------
\makepagestyle{seminar}
\makeoddhead{seminar}{\footnotesize LLM Theory Seminar}{ }{\footnotesize Transformer Expressivity}
\makeoddfoot{seminar}{}{\footnotesize\thepage}{}
\makeheadrule{seminar}{\textwidth}{0.35pt}
\pagestyle{seminar}
\aliaspagestyle{chapter}{seminar}

% ---------- chapter / section hierarchy ----------
\makechapterstyle{seminarnotes}{
  \setlength{\beforechapskip}{8pt}
  \setlength{\midchapskip}{8pt}
  \setlength{\afterchapskip}{22pt}
  \renewcommand{\chapterheadstart}{\vspace*{\beforechapskip}}
  \renewcommand{\chapnamefont}{\normalfont\small\bfseries}
  \renewcommand{\chapnumfont}{\normalfont\bfseries\fontsize{34}{38}\selectfont\color{NoteBlue}}
  \renewcommand{\chaptitlefont}{\normalfont\huge\bfseries\color{black}}
  \renewcommand{\printchaptername}{}
  \renewcommand{\chapternamenum}{}
  \renewcommand{\printchapternum}{\chapnumfont\thechapter}
  \renewcommand{\afterchapternum}{\par\nobreak\color{black}\vskip\midchapskip}
  \renewcommand{\printchaptertitle}[1]{\chaptitlefont ##1}
  \renewcommand{\afterchaptertitle}{\par\nobreak\vskip\afterchapskip}
}
\chapterstyle{seminarnotes}
\setsecheadstyle{\Large\bfseries}
\setsubsecheadstyle{\large\bfseries}
\setsubsubsecheadstyle{\normalsize\bfseries}
\setbeforesecskip{-1.3\baselineskip plus -0.2\baselineskip minus -0.1\baselineskip}
\setaftersecskip{0.55\baselineskip}
\setbeforesubsecskip{-0.9\baselineskip plus -0.15\baselineskip minus -0.1\baselineskip}
\setaftersubsecskip{0.35\baselineskip}

% ---------- notation ----------
\newcommand{\R}{\mathbb{R}}
\newcommand{\N}{\mathbb{N}}
\newcommand{\E}{\mathbb{E}}
\newcommand{\X}{\mathbf{X}}
\newcommand{\Q}{\mathbf{Q}}
\newcommand{\K}{\mathbf{K}}
\newcommand{\V}{\mathbf{V}}
\newcommand{\A}{\mathbf{A}}
\newcommand{\W}{\mathbf{W}}
\newcommand{\one}{\mathbf{1}}
\newcommand{\softmax}{\operatorname{softmax}}
\newcommand{\hardmax}{\operatorname{hardmax}}
\newcommand{\ReLU}{\operatorname{ReLU}}
\newcommand{\Attn}{\operatorname{Attn}}
\newcommand{\FF}{\operatorname{FF}}
\newcommand{\rank}{\operatorname{rank}}
\newcommand{\diag}{\operatorname{diag}}
\newcommand{\norm}[1]{\left\lVert #1\right\rVert}
\newcommand{\abs}[1]{\left\lvert #1\right\rvert}
\newcommand{\inner}[2]{\left\langle #1,#2\right\rangle}
\newcommand{\dd}{\,\mathrm{d}}
\DeclareMathOperator*{\argmax}{arg\,max}
\DeclareMathOperator*{\argmin}{arg\,min}

% ---------- note environments ----------
\tcbset{
  seminarbox/.style={
    enhanced,breakable,sharp corners,boxrule=0pt,frame hidden,
    left=3.2mm,right=3mm,top=1.8mm,bottom=1.8mm,
    before={\par\Needspace{5\baselineskip}},before skip=0.8em,after skip=0.8em,
    fonttitle=\bfseries\small,coltitle=black,
    attach title to upper={\quad},
  }
}
\newtcolorbox{keybox}[1][]{
  seminarbox,colback=NoteBlueLight,
  borderline west={1.8pt}{0pt}{NoteBlue},
  title={핵심#1}
}
\newtcolorbox{paperbox}[1][]{
  seminarbox,colback=NoteGrayLight,
  borderline west={1.8pt}{0pt}{NoteGray},
  title={원 논문 기준#1}
}
\newtcolorbox{suppbox}[1][]{
  seminarbox,colback=NoteWarmLight,
  borderline west={1.8pt}{0pt}{NoteWarm},
  title={보충 유도 --- 논문 밖의 독자용 전개#1}
}
\newtcolorbox{warningbox}[1][]{
  seminarbox,colback=NoteGrayLight,
  borderline west={1.8pt}{0pt}{NoteGray},
  title={주의#1}
}
\newtcolorbox{presentbox}{
  seminarbox,colback=white,
  borderline west={2.2pt}{0pt}{black!70},
  fontupper=\footnotesize,
  top=1.2mm,bottom=1.2mm,before skip=0.5em,after skip=0.5em,
  title={발표에서 이렇게 말하기}
}
\newtcolorbox{presentboxcompact}{
  enhanced,sharp corners,boxrule=0pt,frame hidden,colback=white,
  borderline west={2.2pt}{0pt}{black!70},
  left=3.2mm,right=3mm,top=0.8mm,bottom=0.8mm,
  before={\par},before skip=0.35em,after skip=0.35em,
  fonttitle=\bfseries\small,coltitle=black,
  fontupper=\footnotesize,
  attach title to upper={\quad},
  title={발표에서 이렇게 말하기}
}

\theoremstyle{definition}
\newtheorem{definition}{정의}[chapter]
\newtheorem{proposition}{명제}[chapter]
\newtheorem{theorem}{정리}[chapter]
\newtheorem{lemma}{보조정리}[chapter]

\begin{document}

\begin{titlingpage}
\thispagestyle{empty}
\vspace*{1.2cm}
{\small\bfseries LLM Theory Seminar \hfill Week 4\par}
\vspace{1.4cm}
{\HUGE\bfseries Transformer Expressivity\par}
\vspace{0.55cm}
{\large universal approximation, contextual mapping, positional encoding, depth--width 조건\par}
\vspace{0.9cm}
\rule{\textwidth}{0.7pt}
\vspace{1.1cm}

\begin{minipage}{0.86\textwidth}
\small
이번 주의 질문은 ``Transformer가 실제로 잘 학습되는가?''가 아니라 더 먼저 오는 질문이다.
\textbf{주어진 architecture가 어떤 함수들을 표현할 수 있는가?} Yun et al.은 positional encoding이 없는
Transformer가 연속 permutation-equivariant sequence-to-sequence 함수를 보편적으로 근사하고,
trainable positional encoding을 넣으면 고정 길이의 임의의 연속 sequence-to-sequence 함수까지 근사할 수 있음을 보였다.
증명의 중심은 self-attention이 입력 전체의 문맥을 구별하는 \textbf{contextual mapping}을 만들고,
token-wise feed-forward layer가 그 식별자를 원하는 출력으로 변환한다는 구조다. 이어서 attention matrix 자체의
표현력과 head dimension의 low-rank bottleneck을 분석해, ``fixed width로 universal''이라는 문장이 실제로 무엇을 뜻하고
무엇을 뜻하지 않는지 분리한다.
\end{minipage}

\vfill
{\small
\textbf{Primary readings}\\[0.25em]
Chulhee Yun, Srinadh Bhojanapalli, Ankit Singh Rawat, Sashank J. Reddi, Sanjiv Kumar,
\textbf{Are Transformers universal approximators of sequence-to-sequence functions?}, ICLR 2020.\\
Srinadh Bhojanapalli, Chulhee Yun, Ankit Singh Rawat, Sashank J. Reddi, Sanjiv Kumar,
\textbf{Low-Rank Bottleneck in Multi-head Attention Models}, ICML 2020.
\par}
\vspace{0.8cm}
{\small September 2026\par}
\end{titlingpage}

\tableofcontents*
\clearpage

% =====================================================================
\chapter{먼저 맞춰 둘 기호와 전체 그림}
% =====================================================================

\section{논문마다 다른 행렬 convention}

Transformer 이론 논문에서 가장 먼저 생기는 혼란은 token을 행으로 둘지 열로 둘지다.
현대 구현 설명에서는 보통 $X\in\R^{n\times d}$로 token을 행으로 두지만,
Yun et al.과 Bhojanapalli et al.은
\[
X\in\R^{d\times n}
\]
처럼 token을 \textbf{열}로 둔다. 이 노트는 원 논문의 증명을 따라가기 위해 후자를 기본으로 쓴다.

\begin{center}
\small
\begin{tabularx}{\textwidth}{>{\bfseries\raggedright\arraybackslash}p{0.20\textwidth}Y Y}
\toprule
개념 & 이 노트 / Yun et al. & 흔한 구현 표기 \\
\midrule
입력 sequence & $X\in\R^{d\times n}$ & $X_{\mathrm{row}}\in\R^{n\times d}$ \\
token 위치 & column $X_{:,j}$ & row $X_{j,:}$ \\
permutation & $XP$ & $PX_{\mathrm{row}}$ \\
query/key projection & $W_QX,\,W_KX$ & $X_{\mathrm{row}}W_Q$ 등 \\
attention score & $(W_KX)^\top(W_QX)$ & $QK^\top$ \\
softmax & column-wise & row-wise가 흔함 \\
sequence length & $n$ & $n$ \\
embedding dimension & $d$ & $d$ \\
head size & $m$ 또는 $d_h$ & $d_h$ \\
FFN hidden width & $r$ & $d_{\mathrm{ff}}$ \\
\bottomrule
\end{tabularx}
\end{center}

\begin{warningbox}
행/열 convention을 섞으면 permutation-equivariance 식이 $T(XP)=T(X)P$인지
$T(PX)=PT(X)$인지부터 바뀐다. 둘은 같은 내용이지만 같은 식은 아니다.
이번 노트에서는 원 논문과 동일하게 \emph{token = column}을 쓴다.
\end{warningbox}

\section{이번 주의 세 층위}

세 가지 질문을 분리해야 한다.

\begin{enumerate}
  \item \textbf{한 attention layer의 표현력}: 주어진 입력 $X$에 대해 어떤 attention matrix를 만들 수 있는가?
  \item \textbf{깊은 Transformer의 함수 표현력}: 여러 attention/FFN block을 쌓으면 어떤 seq2seq 함수를 근사할 수 있는가?
  \item \textbf{학습 가능성}: 그 함수를 표현하는 parameter를 SGD가 실제로 찾는가?
\end{enumerate}

이번 주의 universal approximation theorem은 2번에 대한 결과다. 3번을 보장하지 않는다.

\begin{keybox}
\textbf{Expressivity}는 ``좋은 parameter가 존재하는가?''의 문제다.
\textbf{Optimization}은 ``그 parameter를 학습 과정이 찾는가?''의 문제다.
\textbf{Generalization}은 ``찾은 함수가 보지 않은 데이터에서도 맞는가?''의 문제다.
이 셋은 서로 다른 theorem이 필요하다.
\end{keybox}

\section{이번 주의 결론을 먼저 한 줄로}

Yun et al.의 핵심은 다음 두 정리다.

\[
\text{No positional encoding}
\quad\Longrightarrow\quad
\text{continuous permutation-equivariant functions의 UAT},
\]

\[
\begin{aligned}
\text{Trainable positional encoding}
&\quad\Longrightarrow\quad\text{arbitrary continuous seq2seq UAT}\\
&\hspace{4.5em}\text{(on a compact domain).}
\end{aligned}
\]

두 경우 모두 놀라운 점은 attention head 수, head size, FFN width를 크게 만들 필요가 없다는 것이다.
원 논문의 constructive theorem은
\[
(h,m,r)=(2,1,4)
\]
라는 고정 폭으로도 충분하다고 말한다. 대신 필요한 \textbf{depth}는 매우 커질 수 있다.

\begin{presentbox}
``이번 주에는 Transformer를 계산 복잡도나 Turing machine 관점에서 보지 않습니다. 먼저 고정 길이의 연속 함수 근사 문제로 보고,
attention이 문맥을 구별하는 식별자를 만들 수 있기 때문에 좁은 폭의 Transformer도 충분히 깊으면 universal하다는 결과를 봅니다.
다만 이는 존재 정리이지 학습 가능성 정리는 아닙니다.''
\end{presentbox}

% =====================================================================
\chapter{Transformer block의 수학과 permutation symmetry}
% =====================================================================

\section{원 논문 형태의 self-attention}

입력은
\[
X=[x_1,\dots,x_n]\in\R^{d\times n}
\]
이다. $h$개의 attention head를 쓰고 각 head의 크기를 $m$이라 하자.
$i$번째 head에 대해
\[
W_Q^i,W_K^i,W_V^i\in\R^{m\times d},
\qquad
W_O^i\in\R^{d\times m}.
\]
Yun et al.은 residual을 포함한 attention layer를
\[
\Attn(X)
=
X+\sum_{i=1}^h
W_O^iW_V^iX\,
\softmax\!\left((W_K^iX)^\top W_Q^iX\right)
\]
로 쓴다. softmax는 각 column에 적용한다.

현대 구현에서 흔한 $1/\sqrt m$ scaling은 이 표현력 정리의 본질이 아니다. 고정된 양의 scale은 projection matrix에 흡수할 수 있기 때문이다.

\section{token-wise FFN}

attention 결과를 $Z=\Attn(X)$라 쓰면 FFN은
\[
\FF(Z)
=
Z+W_2\ReLU(W_1Z+b_1\one_n^\top)+b_2\one_n^\top,
\]
여기서
\[
W_1\in\R^{r\times d},
\qquad
W_2\in\R^{d\times r}.
\]
중요한 점은 모든 column에 \textbf{같은} $W_1,W_2$가 적용된다는 것이다.
즉 FFN 자체는 token끼리 정보를 섞지 않는다.

\section{왜 positional encoding이 없으면 permutation equivariant인가}

$P\in\R^{n\times n}$를 permutation matrix라 하자. 입력 token 순서를 바꾸면 $X\mapsto XP$이다.

먼저 query와 key는
\[
W_Q(XP)=(W_QX)P,
\qquad
W_K(XP)=(W_KX)P.
\]
따라서 score matrix는
\begin{align*}
S(XP)
&=(W_KXP)^\top(W_QXP)\\
&=P^\top(W_KX)^\top(W_QX)P\\
&=P^\top S(X)P.
\end{align*}

column-wise softmax는 같은 permutation에 대해
\[
\softmax(P^\top SP)=P^\top\softmax(S)P
\]
를 만족한다. 따라서 한 head의 value aggregation은
\begin{align*}
W_V(XP)\softmax(S(XP))
&=W_VXP\,P^\top\softmax(S(X))P\\
&=W_VX\softmax(S(X))P.
\end{align*}
결국 residual까지 포함하면
\[
\Attn(XP)=\Attn(X)P.
\]
FFN도 같은 parameter를 column-wise로 적용하므로
\[
\FF(ZP)=\FF(Z)P.
\]
따라서 block의 합성 역시
\[
T(XP)=T(X)P.
\]

\begin{proposition}[Permutation equivariance]
Positional encoding이 없는 Transformer block과 그 finite composition은 permutation equivariant하다.
즉 모든 permutation matrix $P$에 대해
\[
T(XP)=T(X)P.
\]
\end{proposition}

\begin{suppbox}
위 유도는 원 논문의 Claim 1을 발표용으로 한 줄씩 풀어쓴 것이다. 핵심은 score가 $P^\top SP$로 conjugation되고,
value 쪽의 $P$와 softmax 쪽의 $P^\top$가 곱해져 사라진다는 것이다.
\end{suppbox}

\section{이 symmetry가 의미하는 제한}

예를 들어 positional information이 전혀 없는 상태에서 두 token을 swap하면 출력도 똑같이 swap되어야 한다.
따라서 ``첫 번째 위치만 특별히 0으로 보내라'' 같은 함수는 표현할 수 없다.

즉 UAT를 말하려면 target function class도 이 symmetry를 가져야 한다.
이것이 Yun et al.의 Theorem 2에서 target을 permutation-equivariant 함수로 제한하는 이유다.

\begin{presentbox}
``Positional encoding이 없다는 것은 단순히 위치 정보를 조금 덜 쓴다는 뜻이 아닙니다. Architecture 자체가 permutation equivariance라는 정확한 symmetry를 갖게 됩니다. 그래서 UAT의 target class도 먼저 이 symmetry 안으로 제한해야 합니다.''
\end{presentbox}

% =====================================================================
\chapter{Universal Approximation을 정확히 읽기}
% =====================================================================

\section{고전적 UAT와 이번 정리의 차이}

고전적인 universal approximation theorem은 대략
\[
f:K\subset\R^D\to\R^M
\]
인 연속 함수를 neural network가 임의의 오차까지 근사할 수 있다는 존재 정리다.
Transformer에서는 입력과 출력이 행렬
\[
X\in\R^{d\times n}
\]
이므로 전체 공간을 펼치면 $D=dn$이다. 하지만 Transformer는 token 간 parameter sharing과 attention 구조를 갖기 때문에
일반 MLP의 UAT를 그대로 적용하는 것으로는 architecture의 역할을 설명하지 못한다.

\section{\texorpdfstring{$L^p$}{Lp} 함수 거리}

Yun et al.은 두 함수 $f_1,f_2:\R^{d\times n}\to\R^{d\times n}$ 사이의 거리를
\[
d_p(f_1,f_2)
=
\left(\int
\norm{f_1(X)-f_2(X)}_p^p\dd X\right)^{1/p}
\]
로 둔다. 여기서 matrix norm은 entry-wise $\ell_p$ norm이다.

따라서 ``universal''은
\[
\forall \varepsilon>0,
\quad
\exists\ \theta
\quad\text{s.t.}\quad
 d_p(f,T_\theta)<\varepsilon
\]
라는 뜻이다.

\begin{warningbox}
이 결과는 uniform norm $\sup_X\norm{f(X)-T(X)}$에 대한 theorem이 아니다.
또한 approximation error가 작은 parameter가 \emph{존재}한다고 말할 뿐,
그 parameter를 polynomial time에 찾을 수 있다는 뜻도 아니다.
\end{warningbox}

\section{함수 class 정의}

$\mathcal F_{\mathrm{PE}}$를 compact support를 갖는 연속 permutation-equivariant 함수들의 class라 하자.
즉
\[
f(XP)=f(X)P
\]
이고, 충분히 큰 compact set 밖에서는 $f(X)=0$이다.

한편 positional encoding을 넣은 경우에는 compact domain $K\subset\R^{d\times n}$ 위의 임의의 연속 함수 class를
$\mathcal F_{\mathrm{CD}}$라 쓴다.

\section{Theorem 2: positional encoding 없이도 fixed-width UAT}

\begin{theorem}[Yun et al., 2020, Theorem 2]
$1\le p<\infty$이고 $\varepsilon>0$라 하자.
모든 $f\in\mathcal F_{\mathrm{PE}}$에 대해
\[
g\in\mathcal T_{2,1,4}
\]
가 존재하여
\[
d_p(f,g)\le\varepsilon
\]
를 만족한다.
\end{theorem}

$\mathcal T_{h,m,r}$는 head 수 $h$, 각 head 크기 $m$, FFN hidden width $r$인 Transformer block들을 임의의 깊이로 합성한 함수 class다.
그러므로 정리는
\[
h=2,\qquad m=1,\qquad r=4
\]
라는 \textbf{고정 폭}만으로 충분하다고 말한다.

\begin{keybox}
여기서 ``fixed width''는 \textbf{fixed depth}가 아니다.
$h,m,r$는 $n,d$와 무관한 상수지만, 원하는 오차가 작아질수록 construction의 depth는 커진다.
\end{keybox}

\section{왜 이 결과가 처음 보면 이상한가}

두 가지 제약이 동시에 있다.

\begin{itemize}
  \item FFN은 각 token에 똑같은 함수만 적용한다.
  \item attention은 pairwise dot product와 shared projection만 이용한다.
\end{itemize}

그런데도 모든 연속 equivariant seq2seq function을 근사할 수 있다.
이 역설을 푸는 핵심이 contextual mapping이다.

\begin{presentbox}
``Universal approximation에서 놀라운 숫자는 2, 1, 4입니다. 두 개의 scalar attention head와 폭 4의 FFN이면 충분합니다. 그러나 width가 작다는 말과 network가 작다는 말은 다릅니다. 이 construction은 depth를 크게 써서 표현력을 얻습니다.''
\end{presentbox}

% =====================================================================
\chapter{증명의 중심: contextual mapping}
% =====================================================================

\section{왜 token-wise FFN만으로는 부족한가}

두 문장을 생각하자.
\[
L^{(1)}=[v_{\mathrm{I}},v_{\mathrm{am}},v_{\mathrm{happy}}],
\qquad
L^{(2)}=[v_{\mathrm{I}},v_{\mathrm{am}},v_{\mathrm{Bob}}].
\]
첫 token ``I''의 raw embedding은 두 문장에서 같다. token-wise FFN만 쓰면 같은 입력은 같은 출력으로 보내기 때문에
두 문맥의 ``I''를 다르게 처리할 수 없다.

따라서 먼저 전체 sequence를 보고
\[
(\text{sequence},\text{position/token})
\]
조합마다 서로 다른 scalar 또는 vector code를 만들어야 한다.

\section{Contextual mapping의 정의}

유한 집합 $\mathcal L\subset\R^{d\times n}$에 대해
\[
q:\mathcal L\to\R^{1\times n}
\]
가 contextual mapping이라는 것은 다음을 뜻한다.

\begin{enumerate}
  \item 같은 sequence $L$ 안의 $n$개 값 $q(L)_1,\dots,q(L)_n$이 모두 다르다.
  \item 서로 다른 context $L,L'$가 permutation 관계가 아니라면, $q(L)$과 $q(L')$에 등장하는 값들도 서로 겹치지 않는다.
\end{enumerate}

즉 각 token에 붙은 scalar가
\[
\text{``나는 어느 sequence의 어느 token인가''}
\]
를 식별하는 ID가 된다.

\begin{keybox}
Self-attention의 역할은 최종 함수를 직접 계산하는 것이 아니라,
\textbf{문맥에 따라 달라지는 고유 식별자}를 만드는 것이다.
그 뒤의 token-wise FFN은 이 identifier를 lookup key처럼 사용하여 원하는 출력값을 생성할 수 있다.
\end{keybox}

\section{UAT proof의 세 단계}

Yun et al.의 Theorem 2 proof는 크게 세 단계다.

\subsection{Step 1: 연속 함수를 grid 위의 piecewise-constant 함수로 근사}

compact support를 편의상 $[0,1]^{d\times n}$ 안에 둔다. grid size $\delta>0$에 대해
\[
\mathcal G_\delta
=
\{0,\delta,2\delta,\dots,1-\delta\}^{d\times n}
\]
를 잡는다.
각 grid point $L$에 대해 cube
\[
S_L
=
\prod_{j=1}^d\prod_{k=1}^n
[L_{jk},L_{jk}+\delta)
\]
를 정의한다.

연속성 때문에 $\delta$가 충분히 작으면 같은 cell 안에서 함수값의 변화가 작다.
따라서 적절한 상수 행렬 $A_L$을 골라
\[
\bar f(X)
=
\sum_{L\in\mathcal G_\delta}
A_L\,\mathbf 1\{X\in S_L\}
\]
로 근사할 수 있고,
\[
d_p(f,\bar f)<\varepsilon/3
\]
가 되도록 할 수 있다.

\begin{suppbox}
연속 함수의 grid approximation은 Transformer 고유의 부분이 아니다.
compact set 위에서 uniform continuity를 사용하면 cell diameter가 충분히 작을 때 함수값 차이가 작아진다.
이번 proof에서 진짜 Transformer-specific한 부분은 Step 2다.
\end{suppbox}

\subsection{Step 2: modified Transformer로 grid function을 구현}

분석을 위해 softmax를 hardmax로, ReLU를 특정 piecewise-linear activation으로 바꾼 modified Transformer를 생각한다.
그 안에서는 다시 세 작업을 수행한다.

\begin{enumerate}
  \item token-wise FFN으로 입력을 grid point로 \textbf{quantize}한다.
  \item self-attention stack으로 각 grid sequence에 대한 \textbf{contextual mapping}을 만든다.
  \item token-wise FFN으로 contextual ID를 원하는 $A_L$의 각 column으로 \textbf{value map}한다.
\end{enumerate}

두 번째 단계가 핵심이다.

\subsection{Step 3: hardmax/piecewise activation을 실제 softmax/ReLU Transformer로 근사}

hardmax는 temperature를 높인 softmax로 근사할 수 있다. 예를 들어 최대값이 유일할 때
\[
\softmax(\lambda z)
\xrightarrow[\lambda\to\infty]{}
e_{\argmax z}.
\]
또한 필요한 piecewise-linear scalar map은 작은 ReLU subnetwork로 근사/구현할 수 있다.
이 때문에 modified network를 실제 Transformer $\mathcal T_{2,1,4}$로 옮길 수 있다.

마지막에는 triangle inequality를 사용한다.
\begin{align*}
d_p(f,g)
&\le d_p(f,\bar f)
+d_p(\bar f,\bar g)
+d_p(\bar g,g)\\
&\le \frac{2\varepsilon}{3}+O(\delta^{d/p}).
\end{align*}
$\delta$를 충분히 작게 잡으면 최종 오차가 $\varepsilon$보다 작아진다.

\section{Attention이 contextual ID를 만든다는 의미}

원 논문의 Lemma 6은 grid의 거의 모든 입력에 대해 $h=2,m=1$ attention layer들을 여러 개 쌓아 contextual mapping을 구성한다.
핵심 idea는 한 번의 attention으로 모든 문맥을 완전히 encode한다기보다,
여러 층의 \textbf{selective shift}를 통해 grid point의 정보를 순차적으로 scalar code에 누적하는 것이다.

한 column의 scalar label을
\[
l_j=u^\top L_{:,j}
\]
로 두고, attention operation으로 특정 label만 골라 현재 global max/min 차이에 비례한 shift를 준다.
반복하면 각 token label이 전체 sequence에 의존하는 서로 다른 값
\[
\tilde l_1<\tilde l_2<\cdots<\tilde l_n
\]
로 바뀐다. 더 중요한 점은 서로 permutation 관계가 아닌 두 sequence의 label set도 서로 분리되도록 construction할 수 있다는 것이다.

\begin{warningbox}
여기서 ``attention이 한 층만으로 arbitrary context를 완벽히 encode한다''고 말하면 틀린다.
Yun et al.의 UAT construction은 \textbf{attention layer의 깊은 합성}을 사용한다.
\end{warningbox}

\begin{presentbox}
``증명은 세 단계로 보면 됩니다. 먼저 연속 함수를 grid lookup table로 바꾸고, attention이 각 context-token 쌍에 unique ID를 붙이고,
FFN이 그 ID를 원하는 출력으로 lookup합니다. 마지막으로 proof 편의를 위해 쓴 hardmax를 softmax로 다시 근사합니다.''
\end{presentbox}

% =====================================================================
\chapter{Positional encoding이 symmetry를 어떻게 깨는가}
% =====================================================================

\section{왜 위치 정보가 필요할까}

positional encoding이 없으면 항상
\[
T(XP)=T(X)P
\]
여야 한다. 따라서 arbitrary seq2seq function의 UAT는 불가능하다.
예를 들어
\[
f(X)_{:,1}=0,
\qquad
f(X)_{:,2}=X_{:,2}
\]
처럼 특정 absolute position을 특별 취급하는 함수는 permutation equivariant하지 않다.

\section{trainable positional encoding}

Yun et al.은
\[
E\in\R^{d\times n}
\]
를 trainable positional encoding으로 두고
\[
g_P(X)=g(X+E),
\qquad
g\in\mathcal T_{h,m,r}
\]
인 함수 class를 생각한다.

이제 column $j$에는 입력 token뿐 아니라 $E_{:,j}$라는 위치별 offset이 더해진다.
따라서 token을 permutation해도 위치 offset이 같이 permutation되는 것이 아니므로 원래의 symmetry가 깨진다.

\section{Theorem 3}

\begin{theorem}[Yun et al., 2020, Theorem 3]
$1\le p<\infty$와 $\varepsilon>0$를 고정하자.
compact domain 위의 임의의 연속 함수
\[
f:K\subset\R^{d\times n}\to\R^{d\times n}
\]
에 대해 trainable positional encoding을 갖는
\[
g\in\mathcal T^{P}_{2,1,4}
\]
가 존재하여
\[
d_p(f,g)\le\varepsilon
\]
를 만족한다.
\end{theorem}

즉 positional encoding을 통해 permutation-equivariance라는 structural restriction을 제거하면,
고정 길이 $n$의 sequence space에서는 arbitrary continuous seq2seq function을 근사할 수 있다.

\section{``positional encoding이 있으면 universal''의 정확한 범위}

다음 세 제한을 반드시 기억해야 한다.

\begin{enumerate}
  \item \textbf{고정된 sequence length $n$}: theorem은 variable-length generalization theorem이 아니다.
  \item \textbf{compact domain}: 임의의 무한 domain에서 uniform하게 모든 continuous function을 근사한다는 뜻이 아니다.
  \item \textbf{trainable PE}: 원 theorem의 construction은 trainable positional encoding을 사용한다.
\end{enumerate}

따라서 이 theorem에서 바로
\[
\text{``sinusoidal PE를 쓰는 실제 LLM도 모든 길이에 대해 universal''}
\]
을 결론내릴 수 없다.

\begin{warningbox}
UAT와 length generalization을 연결하지 말 것. 고정 $n$의 함수 공간에서의 universality와,
훈련보다 긴 sequence에서 algorithm을 계속 수행하는 능력은 다른 문제다. 후자는 뒤 주차의 computational theory / length-generalization 문제에 가깝다.
\end{warningbox}

\section{PE는 단순한 부가 정보가 아니라 symmetry-breaking device}

수학적으로 positional encoding의 중요한 역할은
\[
\text{equivariant function class}
\quad\to\quad
\text{general function class}
\]
로 hypothesis class의 symmetry를 바꾸는 것이다.

이 관점은 image model의 translation equivariance, graph network의 permutation invariance에서도 반복된다.
Architecture가 어떤 symmetry를 강제로 보존하는지, positional/node/coordinate feature가 그 symmetry를 어디까지 깨는지를 따지는 것이 표현력 분석의 기본 패턴이다.

\begin{presentbox}
``Positional encoding의 이론적 역할은 위치를 알려주는 것보다 더 정확히 말하면 symmetry를 깨는 것입니다. PE가 없으면 모델 class 전체가 permutation equivariant라 arbitrary function을 표현할 수 없고, trainable PE를 더하면 그 제약이 사라져 fixed-length compact-domain UAT가 됩니다.''
\end{presentbox}

% =====================================================================
\chapter{Attention matrix 자체의 표현력과 low-rank bottleneck}
% =====================================================================

\section{질문을 한 layer 수준으로 내리기}

이제 깊은 Transformer 전체가 아니라 한 attention head를 본다.
$X\in\R^{d\times n}$에 대해 query/key dimension을 $d_k$라 하면
\[
Q=W_QX\in\R^{d_k\times n},
\qquad
K=W_KX\in\R^{d_k\times n}.
\]
attention probability matrix를
\[
P
=
\softmax\!\left(\frac{K^\top Q}{\sqrt{d_k}}\right)
\in\R^{n\times n}
\]
라 하자.

질문은 다음과 같다.

\begin{quote}
고정된 입력 $X$와 원하는 positive column-stochastic matrix $P$가 주어졌을 때,
적절한 $W_Q,W_K$를 골라 정확히 그 $P$를 만들 수 있는가?
\end{quote}

\section{softmax 역문제}

$P$의 모든 원소가 양수이고 각 column 합이 1이라 하자.
만약 score matrix $S$를
\[
S_{ij}=\log P_{ij}
\]
로 두면 각 column에 대해
\[
\softmax(S)_{ij}
=
\frac{e^{\log P_{ij}}}{\sum_{r=1}^n e^{\log P_{rj}}}
=
\frac{P_{ij}}{1}
=P_{ij}.
\]
따라서 문제는 결국
\[
K^\top Q=\sqrt{d_k}\,S
\]
를 만들 수 있는지의 rank/factorization 문제로 바뀐다.

\section{Representation Theorem}

\begin{theorem}[Bhojanapalli et al., 2020, Theorem 1의 핵심]
$X\in\R^{d\times n}$이 full column rank이고 $d\ge n$이라 하자.
query/key projection dimension이 충분히 커서 원하는 $n\times n$ score matrix를 factorize할 수 있으면,
임의의 positive column-stochastic matrix $P$에 대해 적절한 $W_Q,W_K$가 존재하여
\[
\softmax\!\left(\frac{(W_KX)^\top(W_QX)}{\sqrt{d_k}}\right)=P
\]
를 만족한다. 반대로 projection dimension이 sequence length보다 작으면 일반적으로 표현할 수 없는 $(X,P)$가 존재한다.
\end{theorem}

\section{충분한 차원에서의 explicit construction}

$X$가 full column rank이므로 left inverse
\[
X^\dagger=(X^\top X)^{-1}X^\top
\]
가 존재하고
\[
X^\dagger X=I_n.
\]
원하는 score matrix를 $S$라 하자.
$S$를
\[
S=\widetilde K^\top\widetilde Q
\]
로 factorize할 수 있다고 하자. 그 다음
\[
W_K=\widetilde KX^\dagger,
\qquad
W_Q=\widetilde QX^\dagger
\]
로 두면
\begin{align*}
(W_KX)^\top(W_QX)
&=(\widetilde KX^\dagger X)^\top(\widetilde QX^\dagger X)\\
&=\widetilde K^\top\widetilde Q\\
&=S.
\end{align*}
따라서 softmax를 통해 원하는 $P$를 얻는다.

\begin{suppbox}
원 논문은 softmax의 column-wise additive constant 불변성까지 포함해 construction을 쓴다.
발표에서는 ``left inverse로 입력 $X$의 영향을 제거한 뒤 원하는 log-attention matrix를 query/key inner product로 factorize한다''고 설명하면 구조가 가장 잘 보인다.
\end{suppbox}

\section{왜 \texorpdfstring{$d_k<n$}{dk<n}이면 bottleneck이 생기는가}

score matrix는
\[
S=K^\top Q
\]
이므로
\[
\rank(S)\le d_k.
\]
$d_k<n$이면 모든 $n\times n$ matrix를 만들 수 없다. softmax가 nonlinear이긴 하지만,
그 전 단계 logit이 low-rank factorization을 거쳐야 한다는 structural constraint는 남는다.
원 논문은 실제로 특정 $(X,P)$가 어떤 $W_Q,W_K$로도 표현되지 않음을 보인다.

\section{multi-head에서 생기는 직관}

표준 heuristic에서는 embedding dimension $d_{\mathrm{model}}$을 head 수 $h$로 나누어
\[
d_k\approx \frac{d_{\mathrm{model}}}{h}
\]
로 둔다. $h$를 늘리면 각 head의 dimension이 작아진다.
따라서 head 수를 무작정 늘리는 것이 개별 head의 표현력을 항상 키우는 것은 아니다.

\begin{warningbox}
``$d_k<n$이면 Transformer가 universal하지 않다''고 결론내리면 안 된다.
Bhojanapalli et al.의 theorem은 \textbf{한 head가 한 입력에서 arbitrary attention matrix를 정확히 표현하는 능력}에 관한 결과다.
Yun et al.의 UAT는 오히려 $m=1$인 작은 head를 \textbf{깊게 여러 층} 쌓아 universality를 얻는다.
두 결과는 모순이 아니라 분석 대상이 다르다.
\end{warningbox}

\begin{presentbox}
``한 layer만 보면 head dimension이 작을 때 rank bottleneck이 분명히 있습니다. 그런데 전체 network의 universality와 한 layer의 arbitrary attention-matrix expressivity는 다른 개념입니다. Yun의 결과는 작은 head를 깊게 조합해 이 한 층의 제한을 우회합니다.''
\end{presentbox}

% =====================================================================
\chapter{Depth와 Width: fixed-width universal의 실제 비용}
% =====================================================================

\section{폭이 작다는 것과 효율적이라는 것은 다르다}

Theorem 2의 핵심 숫자는 $(h,m,r)=(2,1,4)$다. 그러나 approximation tolerance를 줄일수록 grid resolution $\delta$를 작게 해야 하고,
필요한 block 수는 빠르게 증가한다.

원 논문의 construction은 대략
\[
O\!\left(
 n\frac{(1/\delta)^{dn}}{n!}
\right)
\]
개의 Transformer block을 요구한다. positional encoding을 넣은 general continuous function construction도
\[
O\!\left(n(1/\delta)^{dn}\right)
\]
수준의 깊이를 사용할 수 있다.

따라서 이 theorem은
\[
\text{constant width suffices}
\]
를 보여주지만
\[
\text{polynomial-size efficient representation}
\]
을 보여주는 theorem은 아니다.

\section{왜 curse of dimensionality가 보이는가}

input space의 dimension은 펼치면 $dn$이다. 한 축마다 grid spacing이 $\delta$이면 cell 수가
\[
(1/\delta)^{dn}
\]
으로 증가한다. arbitrary continuous function을 lookup-table식으로 근사하면 이 exponential dependence는 자연스럽다.

즉 universal approximation은 매우 약한 의미의 ``가능성''이다. 실제 언어 함수가 arbitrary worst-case continuous function보다 훨씬 structure가 많기 때문에
실제 모델이 훨씬 적은 parameter로 작동할 수 있는 것이다.

\section{Depth의 역할}

Yun proof에서 depth는 단순한 반복이 아니라 세 가지 계산을 순차적으로 수행하는 resource다.

\begin{enumerate}
  \item 여러 FFN layer가 coordinate별 quantization을 수행한다.
  \item 여러 attention layer가 selective shift를 반복해 context ID를 만든다.
  \item 여러 FFN layer가 많은 ID-to-output mapping을 구현한다.
\end{enumerate}

따라서 ``attention 한 번이 모든 관계를 parallel하게 본다''는 실무적 직관과
``UAT construction이 깊이를 많이 필요로 한다''는 수학적 사실은 동시에 참이다.

\section{width condition을 어떻게 읽을 것인가}

두 paper의 width 관련 문장을 함께 놓으면 다음처럼 구분된다.

\begin{center}
\small
\begin{tabularx}{\textwidth}{>{\bfseries\raggedright\arraybackslash}p{0.24\textwidth}Y Y}
\toprule
질문 & 필요한 width 조건 & 해석 \\
\midrule
깊은 Transformer 전체의 UAT & $h=2,m=1,r=4$로 충분 & depth가 arbitrary하게 커질 수 있음 \\
한 attention head가 arbitrary $P$ 표현 & head dimension이 sequence length 수준이면 충분 & single-layer exact attention expressivity \\
실제 효율적 모델 크기 & 위 theorem들만으로 결정 불가 & data/task structure와 optimization이 필요 \\
\bottomrule
\end{tabularx}
\end{center}

\begin{keybox}
표현력에서 width와 depth는 대체 가능한 resource일 수 있다.
``small width is universal''은 ``width는 중요하지 않다''가 아니라 ``충분한 depth가 주어지면 width를 상수로 고정해도 존재 정리는 성립한다''는 뜻이다.
\end{keybox}

\begin{presentbox}
``Fixed-width UAT의 가장 중요한 caveat는 depth입니다. 폭 4의 FFN과 scalar head만으로 충분하지만, proof는 grid cell 수에 가까운 엄청난 depth를 사용할 수 있습니다. 따라서 universality와 parameter efficiency를 같은 말로 쓰면 안 됩니다.''
\end{presentbox}

% =====================================================================
\chapter{UAT가 말하지 않는 것}
% =====================================================================

\section{1. SGD learnability를 보장하지 않는다}

정리는
\[
\exists\theta:\ d_p(f,T_\theta)<\varepsilon
\]
를 말한다. 하지만 실제 training은
\[
\theta_{t+1}=\theta_t-\eta\nabla_\theta L(\theta_t)
\]
로 움직인다. 존재하는 좋은 $\theta$가 optimization landscape에서 접근 가능한지는 별도 문제다.

지난 주 NTK 결과는 optimization dynamics의 한 regime을 다루었지만, 이번 UAT와 논리적으로 독립이다.

\section{2. sample complexity를 보장하지 않는다}

function class가 universal하면 오히려 너무 많은 함수를 표현할 수 있다.
유한한 data만 보고 target을 식별하려면 capacity control, norm, margin, prior, algorithmic bias 같은 추가 구조가 필요하다.

\section{3. length generalization을 보장하지 않는다}

Theorem 3는 고정된 $n$에 대한 함수 공간의 결과다.
훈련 시 길이 $n\le N$만 보고 $n\gg N$에서도 같은 알고리즘을 수행할 수 있는지는 전혀 다른 문제다.

\section{4. computational complexity를 보장하지 않는다}

어떤 함수를 표현할 수 있다는 것과 그 함수를 깊이/precision/시간 제한 아래 계산할 수 있다는 것은 다르다.
Turing completeness, circuit class $TC^0$, formal-language recognition 같은 결과는 다음 computational theory 주제에서 다루는 것이 맞다.

\section{5. attention weight의 해석 가능성을 보장하지 않는다}

attention matrix를 다양하게 표현할 수 있다는 결과는 attention weight가 유일한 explanation이라는 뜻이 아니다.
value projection과 residual, 여러 head/layer의 조합 때문에 최종 output과 attention weight의 관계는 별도 분석이 필요하다.

\begin{warningbox}
발표에서 가장 위험한 문장은 ``Transformer는 universal하므로 어떤 문제든 풀 수 있다''이다.
정확한 문장은 ``고정 길이의 compact continuous seq2seq function class에 대해 충분한 depth와 적절한 positional encoding을 허용하면 arbitrary accuracy의 approximation parameter가 존재한다''이다.
\end{warningbox}

\section{Expressivity, computability, learnability, generalization}

\begin{center}
\small
\begin{tabularx}{\textwidth}{>{\bfseries\raggedright\arraybackslash}p{0.19\textwidth}Y Y}
\toprule
개념 & 질문 & 대표 수학적 형태 \\
\midrule
Expressivity & 표현 가능한가? & $\exists\theta:\ \norm{T_\theta-f}<\varepsilon$ \\
Computability & 제한된 resource로 계산 가능한가? & depth/precision/circuit/language class \\
Learnability & data와 optimization으로 찾을 수 있는가? & sample/optimization complexity \\
Generalization & unseen input에서도 맞는가? & population risk, OOD/length behavior \\
\bottomrule
\end{tabularx}
\end{center}

\begin{presentbox}
``이번 주 theorem은 architecture의 가능 영역을 정해줍니다. 하지만 가능 영역 안에서 SGD가 무엇을 찾는지, 데이터 몇 개가 필요한지, 더 긴 sequence에서도 되는지는 하나도 자동으로 따라오지 않습니다. 그 분리가 이후 이론 세미나의 출발점입니다.''
\end{presentbox}

% =====================================================================
\chapter{발표 준비용 압축 정리}
% =====================================================================

\section{칠판에 반드시 남길 여덟 식}

\textbf{1. Transformer input convention}
\[
X\in\R^{d\times n}.
\]

\textbf{2. Self-attention score}
\[
S=(W_KX)^\top(W_QX).
\]

\textbf{3. Permutation equivariance}
\[
T(XP)=T(X)P.
\]

\textbf{4. UAT without PE}
\[
\forall f\in\mathcal F_{\mathrm{PE}},\ \forall\varepsilon>0,
\quad
\exists g\in\mathcal T_{2,1,4}:
 d_p(f,g)\le\varepsilon.
\]

\textbf{5. Positional encoding}
\[
g_P(X)=g(X+E).
\]

\textbf{6. UAT with trainable PE}
\[
\forall f\in\mathcal F_{\mathrm{CD}},\ \forall\varepsilon>0,
\quad
\exists g\in\mathcal T^P_{2,1,4}:
 d_p(f,g)\le\varepsilon.
\]

\textbf{7. Attention-logit rank bound}
\[
\rank(K^\top Q)\le d_k.
\]

\textbf{8. Fixed width does not mean fixed size}
\[
\text{width}=O(1),
\qquad
\text{depth can grow as }(1/\delta)^{dn}.
\]

\section{발표 순서 추천}

\begin{enumerate}
  \item Transformer block 수식을 한 번만 정확히 정의한다.
  \item PE가 없을 때 permutation equivariance를 직접 4줄로 증명한다.
  \item UAT의 target class와 $d_p$를 정의한다.
  \item Yun Theorem 2의 $(2,1,4)$ result를 먼저 보여준다.
  \item proof를 quantization $\to$ contextual mapping $\to$ value mapping의 세 단계로 설명한다.
  \item trainable PE가 symmetry를 깨고 Theorem 3으로 이어짐을 보여준다.
  \item Bhojanapalli의 rank theorem으로 한 layer의 attention expressivity를 보완한다.
  \item 마지막에 fixed-width universality와 huge depth, learnability/length generalization의 차이를 분명히 한다.
\end{enumerate}

\section{자주 나오는 질문}

\textbf{Q1. PE가 없으면 Transformer가 순서를 전혀 못 보나?}

absolute position을 구분하지 못하고 permutation equivariant하다. 다만 token content 자체에 순서 정보가 이미 encode되어 있다면 그 정보는 사용할 수 있다.
순수 architecture symmetry에 대한 진술이다.

\textbf{Q2. 두 head, head size 1이면 실제 Transformer도 충분하다는 뜻인가?}

표현 가능성의 존재 정리에서는 충분하다. 그러나 필요한 depth가 매우 클 수 있어 practical efficiency를 의미하지 않는다.

\textbf{Q3. contextual mapping은 그냥 positional encoding인가?}

아니다. contextual mapping은 같은 token도 전체 sequence가 달라지면 다른 ID를 갖게 하는 \emph{content-dependent} mapping이다.
PE는 position-dependent offset으로 symmetry를 깨는 장치다.

\textbf{Q4. 한 head에서 $d_k<n$이면 왜 UAT와 모순이 아닌가?}

single-head single-layer가 arbitrary attention matrix를 정확히 만드는 능력과, 깊은 network가 arbitrary function을 근사하는 능력은 다른 property다.
깊은 합성은 한 층의 rank restriction을 우회할 수 있다.

\textbf{Q5. UAT와 Turing completeness는 같은 말인가?}

아니다. UAT는 연속 공간의 함수 근사 문제이고, Turing completeness는 discrete computation과 unbounded computation model의 문제다.
서로 하나가 다른 하나를 자동으로 함의하지 않는다.

\textbf{Q6. positional encoding이 sinusoidal이어도 Theorem 3가 그대로 성립하나?}

원 theorem은 trainable $E$를 허용한다. fixed sinusoidal PE에 대해 동일한 statement를 그대로 인용하면 안 된다.

\section{흔한 수학적 실수 체크리스트}

\begin{itemize}
  \item row-token / column-token convention을 섞지 않았는가?
  \item PE가 없는 Transformer의 target class를 arbitrary function으로 쓰지 않았는가?
  \item Theorem 2의 metric이 $L^p$ 함수 거리임을 잊지 않았는가?
  \item $(h,m,r)=(2,1,4)$를 fixed depth라고 오해하지 않았는가?
  \item contextual mapping을 한 attention layer의 결과라고 말하지 않았는가?
  \item proof의 hardmax network와 실제 softmax Transformer를 구분했는가?
  \item trainable PE theorem을 sinusoidal/RoPE에 자동 적용하지 않았는가?
  \item single-layer low-rank bottleneck을 전체 deep Transformer의 non-universality로 확대하지 않았는가?
  \item UAT를 optimization/generalization theorem처럼 설명하지 않았는가?
  \item fixed-length UAT를 length generalization claim으로 바꾸지 않았는가?
\end{itemize}

\section{30초 결론}

Positional encoding이 없는 Transformer는 architecture 자체가 permutation equivariant하므로 그 symmetry를 가진 함수만 표현한다.
그 범위 안에서는 Yun et al.이 두 scalar head와 FFN width 4라는 고정 폭만으로도 arbitrary continuous equivariant seq2seq function을 근사할 수 있음을 보였다.
증명의 핵심은 attention stack이 context-dependent unique identifier를 만들고 token-wise FFN이 이를 원하는 output으로 바꾸는 것이다.
Trainable positional encoding을 추가하면 permutation symmetry가 깨져 고정 길이 compact domain의 arbitrary continuous seq2seq function에 대한 UAT로 확장된다.
그러나 fixed-width universality는 매우 큰 depth를 허용하는 존재 정리이며, SGD learnability, sample efficiency, computational complexity, length generalization을 보장하지 않는다.

\begin{presentbox}
``Transformer UAT의 핵심은 attention이 모든 함수를 직접 계산한다는 것이 아니라 context를 구별할 수 있다는 데 있습니다. Attention이 각 context-token 쌍에 고유한 ID를 만들면, 이후 token-wise FFN은 그 ID를 lookup key처럼 써서 arbitrary output을 만들 수 있습니다. PE는 여기에 absolute position까지 구별할 수 있도록 symmetry를 깨줍니다.''
\end{presentbox}

% =====================================================================
\chapter*{참고문헌 및 원문에서 확인할 위치}
\addcontentsline{toc}{chapter}{참고문헌 및 원문에서 확인할 위치}
% =====================================================================

\textbf{Chulhee Yun, Srinadh Bhojanapalli, Ankit Singh Rawat, Sashank J. Reddi, Sanjiv Kumar.}
``Are Transformers universal approximators of sequence-to-sequence functions?'' ICLR 2020.

\begin{itemize}
  \item Sec. 2: Transformer block 정의와 permutation equivariance Claim 1.
  \item Theorem 2: $\mathcal T_{2,1,4}$의 continuous permutation-equivariant UAT.
  \item Theorem 3: trainable positional encoding을 추가한 arbitrary continuous seq2seq UAT.
  \item Def. 3.1, Lemma 6: contextual mapping과 self-attention construction.
  \item Sec. 3.2--4: quantization, contextual mapping, value mapping의 proof skeleton.
  \item Sec. 4.4: fixed width 대신 매우 큰 depth가 필요한 construction의 size discussion.
\end{itemize}
\url{https://arxiv.org/abs/1912.10077}

\vspace{0.8em}
\textbf{Srinadh Bhojanapalli, Chulhee Yun, Ankit Singh Rawat, Sashank J. Reddi, Sanjiv Kumar.}
``Low-Rank Bottleneck in Multi-head Attention Models.'' ICML 2020.

\begin{itemize}
  \item Sec. 2 / Theorem 1: 한 attention head가 arbitrary positive stochastic attention matrix를 표현할 수 있는 dimension 조건.
  \item Proof of Theorem 1: $X^\dagger X=I$를 이용해 input dependence를 제거하고 원하는 score matrix를 구성.
  \item Low-rank case: head projection dimension이 sequence length보다 작을 때 발생하는 representation restriction.
  \item Sec. 3--4: standard multi-head heuristic에서 head 수 증가와 head dimension 감소 사이의 tradeoff.
\end{itemize}
\url{https://proceedings.mlr.press/v119/bhojanapalli20a.html}

\begin{warningbox}[--- 원문과 이 노트의 경계]
Permutation-equivariance의 행렬 유도, softmax inverse를 $\log P$로 보는 설명, grid approximation의 uniform-continuity 직관,
attention score의 rank bound를 통한 bottleneck 설명은 원 theorem을 이해하기 위해 한 줄씩 풀어쓴 보충 전개다.
Yun et al.의 full contextual-mapping proof는 Appendix에서 매우 긴 selective-shift construction을 사용하므로,
이 노트는 그 구조와 theorem statement를 보존하되 모든 index inequality를 복제하지 않았다.
또한 UAT는 고정 길이의 함수 근사 결과이며 Turing completeness나 length generalization 결과로 확장해 해석하지 않는다.
\end{warningbox}

\end{document}
